% =============================================================
% BODY (ES)
% =============================================================

\section{Resolución De La Actividad}
\vspace{2mm}

El paper seleccionado para esta actividad es \textit{Análisis de sentimientos en redes sociales y buscadores online. Enfoque predictivo financiero: un mapeo sistemático de la literatura}. El trabajo analiza cómo el análisis de sentimientos, aplicado sobre redes sociales, motores de búsqueda y otras fuentes digitales, puede complementar los datos financieros tradicionales para mejorar la comprensión y la predicción del comportamiento del mercado.

\vspace{2mm}

A continuación, se responden las cuatro preguntas propuestas por la actividad a partir de la lectura del paper.


\subsection{¿Cuál Era El Problema?}
\vspace{2mm}

El problema central identificado en el paper es la dificultad de predecir con precisión las variaciones de los mercados financieros. Según el artículo, estas variaciones no dependen únicamente de datos históricos o indicadores económicos tradicionales, sino también de múltiples factores externos, como cambios económicos globales, noticias, comportamiento de los inversores y percepciones sociales difundidas en plataformas digitales.

\vspace{2mm}

El paper plantea que muchos operadores financieros consideran que el rendimiento histórico puede utilizarse para anticipar rendimientos futuros. Sin embargo, el trabajo amplía esta mirada al señalar que el sentimiento de los usuarios también cumple un rol importante. En otras palabras, el problema no es solamente técnico o estadístico, sino también informacional: los modelos tradicionales pueden quedar incompletos si no consideran datos no estructurados provenientes de redes sociales, noticias, buscadores online o foros financieros.

\vspace{2mm}

Además, el artículo destaca que los consumidores e inversores son cada vez más difíciles de predecir. Esto vuelve necesario incorporar nuevas fuentes de información capaces de reflejar opiniones, emociones y tendencias en tiempo real. Desde la perspectiva de diseño de software, el problema puede interpretarse como la necesidad de construir soluciones que integren fuentes heterogéneas de datos, procesen lenguaje natural y generen información útil para apoyar decisiones financieras.


\subsection{¿Qué Solución Propone?}
\vspace{2mm}

La solución propuesta no consiste en una única herramienta implementada, sino en un enfoque basado en un mapeo sistemático de la literatura. El objetivo del paper es relevar el estado actual de las técnicas de análisis de sentimientos aplicadas a redes sociales y buscadores online, con orientación a la predicción financiera.

\vspace{2mm}

El artículo propone integrar el análisis de sentimientos con modelos predictivos financieros. Esta integración permite complementar los datos financieros tradicionales con información no estructurada obtenida desde plataformas digitales. El análisis de sentimientos clasifica textos según su polaridad, por ejemplo positiva, negativa o neutra, y esa información puede transformarse en una señal adicional para interpretar tendencias del mercado.

\vspace{2mm}

La solución se apoya en la idea de que redes sociales, noticias web, tendencias de búsqueda y foros pueden reflejar percepciones colectivas sobre empresas, activos o contextos económicos. Al combinar esas señales con modelos predictivos, se busca mejorar la toma de decisiones, anticipar movimientos del mercado y gestionar riesgos con mayor rapidez.

\vspace{2mm}

En el paper se menciona que la precisión de los modelos predictivos puede aumentar cuando se combinan múltiples fuentes de datos, como noticias, redes sociales y otras plataformas en línea. Por lo tanto, la propuesta general es diseñar un enfoque de análisis que no dependa de una única fuente, sino que integre datos financieros tradicionales y datos textuales provenientes del entorno digital.


\subsection{¿Qué Decisiones De Diseño Aparecen?}
\vspace{2mm}

Aunque el paper no presenta el diseño detallado de una aplicación concreta, sí permite identificar varias decisiones de diseño relevantes para una solución basada en análisis de sentimientos aplicado a finanzas.

\vspace{2mm}

En primer lugar, aparece la decisión de utilizar un \textbf{mapeo sistemático de la literatura} como estrategia de investigación. Esto implica definir preguntas de investigación, seleccionar fuentes, establecer cadenas de búsqueda y aplicar criterios de inclusión y exclusión. Esta decisión ordena el análisis y permite justificar los resultados a partir de bibliografía seleccionada.

\vspace{2mm}

En segundo lugar, se decide trabajar con \textbf{fuentes digitales múltiples}. El paper menciona bases y fuentes como SEDICI, IEEE Xplore Digital Library, ACM Digital Library y Science Direct para la revisión bibliográfica. A nivel de una solución predictiva, también se consideran redes sociales, noticias web, tendencias de Google y discusiones en foros. Esta decisión apunta a mejorar la riqueza del conjunto de datos y evitar depender exclusivamente de indicadores financieros clásicos.

\vspace{2mm}

En tercer lugar, aparece la decisión de utilizar \textbf{procesamiento de lenguaje natural} para clasificar textos según sentimiento. Esta decisión es importante porque el sistema debe transformar opiniones, noticias y publicaciones en datos interpretables por modelos predictivos. Desde el punto de vista del diseño, esto supone definir mecanismos de extracción, limpieza, clasificación y representación de información textual.

\vspace{2mm}

También se observa una decisión de diseño relacionada con el \textbf{recorte del dominio}. El paper no analiza el análisis de sentimientos de forma general, sino específicamente aplicado a redes sociales, buscadores online y predicción financiera. Este recorte permite enfocar mejor las preguntas de investigación y los criterios de selección de artículos.

\vspace{2mm}

Finalmente, se identifica la decisión de considerar tanto \textbf{beneficios como riesgos}. El paper no presenta el análisis de sentimientos como una solución automática o perfecta, sino como un enfoque que puede mejorar la toma de decisiones, pero que también trae desafíos vinculados a la complejidad de los modelos, la ambigüedad del lenguaje, la volatilidad del mercado, la transparencia, la privacidad y la capacitación de los analistas financieros.


\subsection{¿Qué Atributos De Calidad Intenta Mejorar?}
\vspace{2mm}

El atributo de calidad más evidente que intenta mejorar el enfoque propuesto es la \textbf{precisión}. El paper busca mostrar cómo la incorporación de datos provenientes de redes sociales, buscadores y noticias puede complementar la información financiera tradicional, permitiendo predicciones más ajustadas sobre tendencias del mercado.

\vspace{2mm}

También se intenta mejorar la \textbf{oportunidad de la información}. Al utilizar fuentes digitales en tiempo real, el enfoque permite detectar cambios de opinión, tendencias emergentes o reacciones públicas antes de que se reflejen completamente en indicadores financieros tradicionales. Esto puede favorecer decisiones más rápidas por parte de inversores, analistas o empresas.

\vspace{2mm}

Otro atributo relevante es la \textbf{capacidad de adaptación}. El paper señala que el análisis de sentimientos puede utilizarse para personalizar estrategias financieras según perfiles de riesgo, preferencias y comportamientos de los inversores. Esto implica que una solución basada en este enfoque debería adaptarse a distintos usuarios y contextos de decisión.

\vspace{2mm}

Además, aparece la \textbf{robustez} como atributo necesario. El artículo reconoce que el lenguaje en redes sociales y noticias puede ser ambiguo, y que los mercados financieros son volátiles y complejos. Por eso, un sistema de este tipo debe ser capaz de distinguir información relevante de ruido, manejar múltiples fuentes y evitar conclusiones débiles basadas en datos subjetivos o inestables.

\vspace{2mm}

También se relaciona con la \textbf{seguridad y privacidad}, especialmente porque en la conclusión se menciona la necesidad de garantizar la privacidad y seguridad de los datos financieros. Este punto es importante porque una solución predictiva financiera puede trabajar con información sensible de usuarios, inversores o instituciones.

\vspace{2mm}

Por último, puede señalarse la \textbf{usabilidad} o comprensibilidad de los resultados. El paper advierte que la complejidad de los modelos debe equilibrarse con la transparencia y la comprensión por parte del usuario. Por lo tanto, no alcanza con que el sistema prediga: también debe explicar o presentar sus resultados de manera clara para que puedan ser usados responsablemente en la toma de decisiones.


\section{Conclusión}
\vspace{2mm}

En síntesis, el paper aborda un problema actual: la dificultad de anticipar comportamientos financieros en contextos altamente variables. La solución propuesta consiste en estudiar e integrar técnicas de análisis de sentimientos con modelos predictivos, aprovechando información no estructurada proveniente de redes sociales, buscadores, noticias y otras fuentes digitales.

\vspace{2mm}

Desde la mirada de diseño de software, el trabajo permite reconocer decisiones vinculadas con la selección de fuentes, el procesamiento de lenguaje natural, la integración de datos, el recorte del dominio y la necesidad de contemplar riesgos. Los principales atributos de calidad asociados son precisión, oportunidad, adaptabilidad, robustez, seguridad, privacidad y claridad en la presentación de resultados.